\section{Coupling Calibration}\label{sec:coupling-calibration}

\subsection{Coupling Framework}\label{ssec:coupling-framework}

The marginal CARMA models are kept fixed. For a weather or renewable factor
\(F\in\{T,S,W\}\), \(dL^F\) denotes the calibrated marginal driver, with its
variance scale included as in \cref{eq:carma-state-space}. The coupled price
model transmits this factor driver into the price driver.
\begin{align}
  dZ_t^F
  &=
  A_F Z_t^F\,dt+e_F\,dL_t^F,
  &
  Y_t^F &= b_F^\top Z_t^F,
  \nonumber \\
  dZ_t^P
  &=
  A_P Z_t^P\,dt
  +
  e_P\left(\lambda_F\,dL_t^F+dL_t^{P,i}\right),
  &
  Y_t^P &= b_P^\top Z_t^P .
  \label{eq:coupling-state-space}
\end{align}
The parameter \(\lambda_F\) acts only on shocks. It does not change the price
drift, the deterministic seasonality, or the marginal CARMA matrices. If the
factor driver is represented as \(dL_t^F=\sigma_F dW_t^F\) in a Gaussian
covariance calculation, the transmitted shock is equivalently
\(\lambda_F\sigma_F e_P dW_t^F\). The paper uses the \(dL\)-driver convention.

The calibration is not performed on the raw deseasonalized residual levels
\(Y^F\) and \(Y^P\). These correlations are only reported as diagnostics. On
hourly intervals, filtered state-space residuals \(R_i^F\) and \(R_i^P\) are
computed from the fixed marginal filters. If \(dL^F\) is independent of the
idiosyncratic price driver \(dL^{P,i}\), their cross-covariance is proportional
to \(\lambda_F\). Define
\[
  D_{FP}
  =
  \frac{1}{n}\sum_{i=1}^{n}
  (R_i^F-\overline R^F)
  (R_i^P-\overline R^P)^\top
\]
as the empirical cross-covariance, and let \(K_{FP}\) denote the theoretical
cross-covariance generated by the same driver scaling when \(\lambda_F=1\).
The scalar estimate is
\begin{equation}
  \lambda_F
  =
  \frac{b_F^\top D_{FP} b_P}
       {b_F^\top K_{FP} b_P}.
  \label{eq:coupling-lambda-estimator}
\end{equation}
The tables also report the projected state-space residual correlation
\(\mathrm{Corr}(b_P^\top R^P,b_F^\top R^F)\), which is the empirical
correlation associated with the covariance used in
\cref{eq:coupling-lambda-estimator}.
The marginal price driver is then decomposed as
\begin{equation}
  dL_i^{P,m}
  =
  dL_i^{P,i}
  +
  \lambda_F\,dL_i^F,
  \qquad
  dL_i^{P,i}
  =
  dL_i^{P,m}
  -
  \lambda_F\,dL_i^F .
  \label{eq:coupling-driver-decomposition}
\end{equation}
After removing \(\lambda_F dL^F\), the remaining correlation with the factor
driver is used as a direct covariance-level diagnostic for the coupled
simulation model.

\subsection{Temperature--Price Coupling}\label{ssec:temperature-price-coupling}

The raw correlation between the deseasonalized temperature residual and the
deseasonalized German log-price residual is \(-0.24\). A negative sign is
consistent with the winter heating channel, as below-seasonal temperatures tend
to increase electricity demand and spot prices. Since the model acts on shocks,
this level correlation is only a diagnostic.

For temperature, a positive \(dL^T\) is a warmer-than-seasonal shock. A negative
\(\lambda_T\) therefore means that warmer-than-seasonal shocks lower price
shocks, or equivalently that colder-than-seasonal shocks raise price shocks.
This is the expected winter sign, although summer cooling demand can create the
opposite local effect. The pooled estimate is
\[
  \lambda_T = -1.4\times 10^{-3}.
\]

\begin{table}[H]
  \centering
  \caption{Temperature--price coupling diagnostics.}
  \label{tab:temperature-price-coupling-diagnostics}
  \begin{tabular}{@{}lr@{}}
    \toprule
    Quantity & Value \\
    \midrule
    Corr.\((Y^P,Y^T)\) & \(-0.24\) \\
    Corr.\((b_P^\top R^P,b_T^\top R^T)\) & \(-0.040\) \\
    Corr.\((dL^{P,m},dL^T)\) & \(-0.042\) \\
    Corr.\((dL^{P,i},dL^T)\) & \(-9.4{\times}10^{-4}\) \\
    Std.\((\lambda_T dL^T)\) / Std.\((dL^{P,i})\) & \(4.1\%\) \\
    \bottomrule
  \end{tabular}
\end{table}

After removing \(\lambda_T dL^T\), the remaining price driver is close
to uncorrelated with the temperature driver, which supports the covariance
decomposition used for simulation. The ratio in the table is therefore close to
the absolute marginal driver correlation for a simple covariance reason. Let
\(X=dL^{P,m}\), \(Y=dL^{P,i}\), and \(Z=\lambda_T dL^T\). Since \(X=Y+Z\) and
\(\operatorname{Cov}(Y,Z)\simeq 0\) after the removal,
\[
  \operatorname{Cov}(X,Z)
  =
  \operatorname{Cov}(Y+Z,Z)
  \simeq
  \operatorname{Var}(Z).
\]
It follows that \(|\operatorname{Corr}(X,Z)|\simeq \operatorname{Std}(Z)/
\operatorname{Std}(X)\). Since the transmitted component is small,
\(\operatorname{Std}(X)\simeq \operatorname{Std}(Y)\), which gives
\[
  |\operatorname{Corr}(dL^{P,m},dL^T)|
  \simeq
  \frac{\operatorname{Std}(\lambda_T dL^T)}
       {\operatorname{Std}(dL^{P,i})}.
\]

\subsection{Solar--Price Coupling}\label{ssec:solar-price-coupling}

For solar, the raw correlation between the deseasonalized latent residual and
the deseasonalized German log-price residual is \(0.063\). The solar factor is
the deseasonalized clear-sky shortfall variable, so a positive residual
corresponds to a larger latent solar shortfall, not to higher solar production.

Because \(Y^S\) increases with latent solar shortfall, a positive solar shock
represents lower effective solar availability. A positive coupling coefficient
would therefore match the merit-order intuition, as lower solar availability raises
residual demand and price shocks. Unlike temperature and wind, the raw residual
correlation is positive, whereas the projected state-space residual correlation
is negative. The coupling estimator is based on the latter covariance, which
explains the negative sign of the estimate.
\[
  \lambda_S = -8.0\times 10^{-6}.
\]

\begin{table}[H]
  \centering
  \caption{Solar--price coupling diagnostics.}
  \label{tab:solar-price-coupling-diagnostics}
  \begin{tabular}{@{}lr@{}}
    \toprule
    Quantity & Value \\
    \midrule
    Corr.\((Y^P,Y^S)\) & \(0.063\) \\
    Corr.\((b_P^\top R^P,b_S^\top R^S)\) & \(-0.0018\) \\
    Corr.\((dL^{P,m},dL^S)\) & \(-1.0{\times}10^{-3}\) \\
    Corr.\((dL^{P,i},dL^S)\) & \(7.8{\times}10^{-4}\) \\
    Std.\((\lambda_S dL^S)\) / Std.\((dL^{P,i})\) & \(0.18\%\) \\
    \bottomrule
  \end{tabular}
\end{table}

The solar driver correlation and the transmitted-shock standard-deviation ratio
are both close to zero. Solar has the smallest ratio among the three factors,
which questions the relevance of a coupled price--solar shock model for this
dataset.

\subsection{Wind--Price Coupling}\label{ssec:wind-price-coupling}

For wind, the raw correlation between the logit residual and the German
log-price residual is \(-0.50\). The wind transform is monotone in the capacity
factor, so a positive wind residual corresponds to above-seasonal wind
production. The negative sign is consistent with the merit-order effect.

For wind, a positive \(dL^W\) is an above-seasonal production shock. The expected
coupling sign is therefore negative, since higher wind production lowers
residual demand. The estimate is
\[
  \lambda_W = -1.7\times 10^{-3}.
\]

\begin{table}[H]
  \centering
  \caption{Wind--price coupling diagnostics.}
  \label{tab:wind-price-coupling-diagnostics}
  \begin{tabular}{@{}lr@{}}
    \toprule
    Quantity & Value \\
    \midrule
    Corr.\((Y^P,Y^W)\) & \(-0.50\) \\
    Corr.\((b_P^\top R^P,b_W^\top R^W)\) & \(-0.040\) \\
    Corr.\((dL^{P,m},dL^W)\) & \(-0.040\) \\
    Corr.\((dL^{P,i},dL^W)\) & \(3.7{\times}10^{-5}\) \\
    Std.\((\lambda_W dL^W)\) / Std.\((dL^{P,i})\) & \(4.0\%\) \\
    \bottomrule
  \end{tabular}
\end{table}

The wind channel is the clearest renewable coupling in this calibration. The
marginal price driver has a small negative wind correlation, and after removing
\(\lambda_W dL^W\) the remaining idiosyncratic driver has no detectable
linear correlation with wind. The transmitted component remains small, with a
standard-deviation ratio close to the temperature one, and this ratio is
consistent with the state-space residual correlation.

\FloatBarrier
